\documentclass{claudeeditorial}

\renewcommand{\DocKicker}{Plantilla}
\renewcommand{\DocTitle}{Título del documento}
\renewcommand{\DocSubtitle}{Subtítulo cálido, claro y editorial para explicar el propósito del artefacto}
\renewcommand{\DocAuthor}{Tu Nombre}
\renewcommand{\DocRole}{Rol o equipo}
\renewcommand{\DocDate}{\today}
\renewcommand{\DocRef}{DOC-001}
\renewcommand{\DocFooter}{claudeeditorial latex}

\hypersetup{
  pdftitle={Título del documento},
  pdfauthor={Tu Nombre}
}

\ClaudeRunningHeader
\setstretch{1.22}

\begin{document}
\BodyCopy

\ClaudeCover

\begin{claudehero}
\ClaudeLead{Escribe aquí el resumen ejecutivo: una idea principal, el contexto suficiente y la promesa del documento. El estilo está pensado para documentos largos que necesitan rigor sin perder calidez.}
\end{claudehero}

\section{Resumen}

Esta sección contiene el cuerpo inicial del documento. Puedes combinar texto normal con \Highlight{realces suaves}, \Tag{etiquetas}, badges, métricas y cajas editoriales.

\begin{tcbraster}[raster columns=3,raster equal height=rows,raster column skip=8pt]
\ClaudeStat{Bloques}{20+}{Componentes reutilizables para informes, propuestas y notas.}
\ClaudeStat{Paleta}{Warm}{Terracota, pergamino, marfil y grises orgánicos.}
\ClaudeStat{Salida}{PDF}{Compilación recomendada con XeLaTeX o LuaLaTeX.}
\end{tcbraster}

\section{Decisión}

\begin{claudedecision}{Criterio principal}
Usa este bloque para fijar una recomendación, una conclusión o una postura clara. El acento terracota debe reservarse para momentos de decisión real.
\end{claudedecision}

\section{Próximos pasos}

\begin{claudechecklist}{Checklist}
\begin{itemize}
\item ajustar metadatos;
\item escribir el resumen ejecutivo;
\item sustituir los ejemplos por contenido real;
\item compilar y revisar saltos de página.
\end{itemize}
\end{claudechecklist}

\ClaudeSignature{\DocAuthor}{\DocRole}

\end{document}
